\documentclass{article}
\usepackage{graphicx}
\usepackage[margin=1.5cm]{geometry}
\usepackage{amsmath}
\usepackage{enumitem}

\setlist{nosep,leftmargin=*}

\begin{document}
\twocolumn

\title{Chapter 4 In-Class Exercise: Sorting Math Students}
\author{Prof. Jordan C. Hanson}
\date{}

\maketitle

\section{Student Data}

\begin{enumerate}

\item A placement test assigns each student a math score from
0 through 5.  Begin with the following list of tuples:

\begin{verbatim}
students = [
    ("Maya", 5),
    ("Noah", 1),
    ("Olivia", 3),
    ("Liam", 4),
    ("Sophia", 2),
    ("Ethan", 0),
    ("Ava", 5),
    ("Lucas", 1),
    ("Emma", 3),
    ("Mateo", 4)
]
\end{verbatim}

Each element of \texttt{students} is a tuple containing

\begin{verbatim}
(student_name, math_score)
\end{verbatim}

For example,

\begin{verbatim}
print(students[0][0])
print(students[0][1])
\end{verbatim}

prints the first student's name and score.

\end{enumerate}


\section{Course Rosters}

\begin{enumerate}

\item Students are assigned according to

\[
\begin{array}{c|c}
\textbf{Score} & \textbf{Course} \\ \hline
4,5 & \text{Calculus} \\
2,3 & \text{Trigonometry} \\
0,1 & \text{Algebra}
\end{array}
\]

Create three course lists.  Each list contains the course
name followed by an initially empty student list:

\begin{verbatim}
algebra = [
    "Algebra", []
]

trigonometry = [
    "Trigonometry", []
]

calculus = [
    "Calculus", []
]
\end{verbatim}

Combine the three lists into one tuple:

\begin{verbatim}
math_courses = (
    algebra,
    trigonometry,
    calculus
)
\end{verbatim}

The final structure will therefore have the form

\begin{verbatim}
(
  ["Algebra", [students...]],
  ["Trigonometry", [students...]],
  ["Calculus", [students...]]
)
\end{verbatim}

\end{enumerate}


\section{Conditional Sorting}

\begin{enumerate}

\item Use conditional logic to place each student into the
correct course.  Begin with the first student:

\begin{verbatim}
name = students[0][0]
score = students[0][1]

if score >= 4:
    calculus[1].append(name)
elif score >= 2:
    trigonometry[1].append(name)
else:
    algebra[1].append(name)
\end{verbatim}

The method

\begin{verbatim}
.append(name)
\end{verbatim}

adds the student's name to the end of the appropriate
student list.

\begin{itemize}

\item Repeat the conditional statement for
\texttt{students[1]} through \texttt{students[9]}.
Do not use a loop yet.

\item For each student, extract the name and score first,
then use one \texttt{if}-\texttt{elif}-\texttt{else}
statement to make the placement.

\item When all ten students have been processed, print

\begin{verbatim}
print(math_courses)
\end{verbatim}

and verify that every student appears exactly once.

\end{itemize}

\end{enumerate}


\section{Check Your Result}

\begin{enumerate}

\item Without changing the original \texttt{students} data,
verify the following:

\begin{itemize}
\item Students with scores 0 or 1 appear only in Algebra.
\item Students with scores 2 or 3 appear only in Trigonometry.
\item Students with scores 4 or 5 appear only in Calculus.
\item The three student lists contain a total of 10 names.
\end{itemize}

\end{enumerate}

\end{document}